% ==================== 第8章 讨论与结论 ====================
\chapter{讨论与结论}
\label{ch:conclusion}

\section{讨论}

本文的实验支持一个分层结论，也决定了论文的主张边界。决策接口是稳健的：融合后验的全动作读出在三个基座上相对受限门分别降低0.749、0.581和0.696的决策损失，随机森林后验降低0.467，融合与随机森林共六个配对区间全部排除零。证据状态是有竞争力的：融合系统在三个基座的准确率、平衡准确率、宏F1和决策损失上均为最高，相对最强基线的点估计更低，但只有Qwen的区间排除零。误差分层提供了更清楚的证据：在真实结果落在指引区间内的82条断言上，三个基座的融合损失都显著低于随机森林，DeepSeek和Qwen的五个种子全部排除零，GLM有四个种子排除零。

这个分层对金融应用有意义。一个总体损失与强结构化基线相当的系统，仍然可以在特定误差层上提供互补信号，并且可以把证据、资格、可靠性和动作完整保留下来。金融决策关心的不只是平均准确率，还包括错误发生在哪个区域、错误成本是否对称、输出是否可复核。本文把这些维度放进同一个契约和同一个损失，因而能够分别报告后验质量、动作空间和证据状态的作用。

因此本文不把全面性能优越作为唯一主张。锚点层的作用应定位为可审计的语义状态和证据资格判定协议，而不是普遍的点预测优势。文本锚点把纯LLM分支的准确率从.445提高到.502，融合后验在区间内样本上显著优于随机森林，但在总体损失上与随机森林的差异不显著。这个结果不削弱框架的价值：监管、审计和研究场景需要的正是一个能够说明“为什么这样判断”的状态，而不是一个只返回标签的黑箱。

外部基准给出了适用边界。FinDVer上，状态条件式获取在DeepSeek和GLM上以约一半的段落数取得高于固定8条的判定准确率，在Qwen3上则低于固定4条。A股回测上，数值桥在DeepSeek上明显改善损失，在GLM上没有触发。两个结果指向同一个条件：证据策略和数值桥的收益取决于基座的抽取能力。框架提供接口，基座决定接口能否被充分利用。

\section{局限}

第一，主基准只覆盖美国公司的季度营业收入指引。A股回测补充了中国市场的年报断言，但两个市场的制度、披露习惯和行业结构不同，跨市场结论仍需更多数据。FinDVer只用于组件检查，不构成端到端前瞻性预测的外部验证。

第二，标签来自XBRL核验或披露文件回退，不是人工金标准。季度口径、重述和单位差异会带来噪声；46条纯文本标签被排除在主结果之外，缩小了测试集。A股回测的数值桥依赖抽取器输出数值字段，GLM上的桥接未触发，Qwen的数值误差较大，说明抽取质量仍是实际部署的主要瓶颈。

第三，类型库扩展只在DeepSeek上完成一轮，GLM和Qwen3没有候选通过验证门控。单轮扩展的测试增益较小，配对区间包含零；更多轮次、跨基座类型迁移和人工审计留待后续工作。

\section{结论}

本文把前瞻性金融断言表示为截止时点上的契约，并在有限结果空间上维护可演化的锚点状态。框架由三层锚点、五类资格判定、揭示日Beta可靠性、验证集门控的类型库扩展、确定性数值桥和全动作损失感知读出组成。实验在227条核验断言、470条A股年报断言和150条FinDVer断言上进行，覆盖三个基座。

结果支持三点结论。第一，全动作读出是稳健的决策改进，在三个基座上显著降低决策损失。第二，融合系统在三个基座的分类和决策指标上均为最高，并在指导区间内样本上显著优于随机森林。第三，类型库扩展和数值桥是条件增益，其效果取决于基座的抽取质量。框架的主要贡献不是某个单点的准确率提升，而是把证据、判断和动作组织成一个可检查、可演化、可复用的决策接口。

\clearpage
\phantomsection
\addcontentsline{toc}{chapter}{参考文献}
{\zihao{5}\setlength{\baselineskip}{16pt}\begin{thebibliography}{28}
\setlength{\bibsep}{0pt}
\setlength{\itemsep}{0pt}
\setlength{\parsep}{0pt}
\setlength{\bibhang}{2em}

\bibitem[Angelopoulos and Bates(2021)]{angelopoulos2021gentle}
Angelopoulos, A. N., and Bates, S., 2021, ``A Gentle Introduction to Conformal Prediction and Distribution-Free Uncertainty Quantification'', arXiv:2107.07511.

\bibitem[Asai et al.(2024)]{asai2024selfrag}
Asai, A., Wu, Z., Wang, Y., Sil, A., and Hajishirzi, H., 2024, ``Self-RAG: Learning to Retrieve, Generate, and Critique through Self-Reflection'', in \emph{International Conference on Learning Representations}.

\bibitem[Brodersen et al.(2010)]{brodersen2010balanced}
Brodersen, K. H., Ong, C. S., Stephan, K. E., and Buhmann, J. M., 2010, ``The Balanced Accuracy and Its Posterior Distribution'', in \emph{Proceedings of the 20th International Conference on Pattern Recognition}, pp. 3121--3124.

\bibitem[Chen et al.(2021)]{chen2021finqa}
Chen, Z., Chen, W., Smiley, C., Shah, S., Borova, I., Langdon, D., Moussa, R., Beane, M., Huang, T.-H., Routledge, B., and Wang, W. Y., 2021, ``FinQA: A Dataset of Numerical Reasoning over Financial Data'', in \emph{Proceedings of the 2021 Conference on Empirical Methods in Natural Language Processing}, pp. 3697--3711.

\bibitem[Chow(1970)]{chow1970optimum}
Chow, C. K., 1970, ``On Optimum Recognition Error and Reject Tradeoff'', \emph{IEEE Transactions on Information Theory}, Vol. 16, No. 1, pp. 41--46.

\bibitem[Elkan(2001)]{elkan2001cost}
Elkan, C., 2001, ``The Foundations of Cost-Sensitive Learning'', in \emph{Proceedings of the 17th International Joint Conference on Artificial Intelligence}, pp. 973--978.

\bibitem[Geifman and El-Yaniv(2017)]{geifman2017selective}
Geifman, Y., and El-Yaniv, R., 2017, ``Selective Classification for Deep Neural Networks'', in \emph{Advances in Neural Information Processing Systems}, Vol. 30.

\bibitem[Gibbs and Cand\`es(2021)]{gibbs2021adaptive}
Gibbs, I., and Cand\`es, E. J., 2021, ``Adaptive Conformal Inference Under Distribution Shift'', in \emph{Advances in Neural Information Processing Systems}, Vol. 34, pp. 1660--1672.

\bibitem[Guo et al.(2017)]{guo2017calibration}
Guo, C., Pleiss, G., Sun, Y., and Weinberger, K. Q., 2017, ``On Calibration of Modern Neural Networks'', in \emph{Proceedings of the 34th International Conference on Machine Learning}, Vol. 70, pp. 1321--1330.

\bibitem[Hirst et al.(2008)]{hirst2008management}
Hirst, D. E., Koonce, L., and Venkataraman, S., 2008, ``Management Earnings Forecasts: A Review and Framework'', \emph{Accounting Horizons}, Vol. 22, No. 3, pp. 315--338.

\bibitem[Islam et al.(2023)]{islam2023financebench}
Islam, P., Kannappan, A., Kiela, D., Qian, R., Scherrer, N., and Vidgen, B., 2023, ``FinanceBench: A New Benchmark for Financial Question Answering'', arXiv:2311.11944.

\bibitem[Jeong et al.(2024)]{jeong2024adaptiverag}
Jeong, S., Baek, J., Cho, S., Hwang, S. J., and Park, J. C., 2024, ``Adaptive-RAG: Learning to Adapt Retrieval-Augmented Large Language Models through Question Complexity'', in \emph{Proceedings of the 2024 Conference of the North American Chapter of the Association for Computational Linguistics}, pp. 7036--7050.

\bibitem[Jiang et al.(2023)]{jiang2023flare}
Jiang, Z., Xu, F., Gao, L., Sun, Z., Liu, Q., Dwivedi-Yu, J., Yang, Y., Callan, J., and Neubig, G., 2023, ``Active Retrieval Augmented Generation'', in \emph{Proceedings of the 2023 Conference on Empirical Methods in Natural Language Processing}, pp. 7969--7992.

\bibitem[Lewis et al.(2020)]{lewis2020rag}
Lewis, P., Perez, E., Piktus, A., Petroni, F., Karpukhin, V., Goyal, N., Kuttler, H., Lewis, M., Yih, W.-T., Rockt\"aschel, T., Riedel, S., and Kiela, D., 2020, ``Retrieval-Augmented Generation for Knowledge-Intensive NLP Tasks'', in \emph{Advances in Neural Information Processing Systems}, Vol. 33, pp. 9459--9474.

\bibitem[L\'opez de Prado(2018)]{lopezdeprado2018advances}
L\'opez de Prado, M., 2018, \emph{Advances in Financial Machine Learning}, Wiley.

\bibitem[Manski(2003)]{manski2003partial}
Manski, C. F., 2003, \emph{Partial Identification of Probability Distributions}, Springer.

\bibitem[Qwen Team(2025)]{qwen2025qwen3}
Qwen Team, 2025, ``Qwen3 Technical Report'', arXiv:2505.09388.

\bibitem[Rangapur et al.(2025)]{rangapur2025finfact}
Rangapur, A., Wang, H., Jian, L., and Shu, K., 2025, ``Fin-Fact: A Benchmark Dataset for Multimodal Financial Fact-Checking and Explanation Generation'', in \emph{Companion Proceedings of the ACM Web Conference 2025}, pp. 785--788.

\bibitem[Reddy et al.(2024)]{reddy2024docfinqa}
Reddy, V., Koncel-Kedziorski, R., Lai, V. D., Krumdick, M., Lovering, C., and Tanner, C., 2024, ``DocFinQA: A Long-Context Financial Reasoning Dataset'', in \emph{Proceedings of the 62nd Annual Meeting of the Association for Computational Linguistics (Volume 2: Short Papers)}, pp. 445--458.

\bibitem[Romano et al.(2020)]{romano2020classification}
Romano, Y., Sesia, M., and Cand\`es, E. J., 2020, ``Classification with Valid and Adaptive Coverage'', in \emph{Advances in Neural Information Processing Systems}, Vol. 33, pp. 3581--3591.

\bibitem[Shinn et al.(2023)]{shinn2023reflexion}
Shinn, N., Cassano, F., Berman, A., Gopinath, A., Narasimhan, K., and Yao, S., 2023, ``Reflexion: Language Agents with Verbal Reinforcement Learning'', in \emph{Advances in Neural Information Processing Systems}, Vol. 36.

\bibitem[Sokolova and Lapalme(2009)]{sokolova2009performance}
Sokolova, M., and Lapalme, G., 2009, ``A Systematic Analysis of Performance Measures for Classification Tasks'', \emph{Information Processing and Management}, Vol. 45, No. 4, pp. 427--437.

\bibitem[Vovk et al.(2005)]{vovk2005algorithmic}
Vovk, V., Gammerman, A., and Shafer, G., 2005, \emph{Algorithmic Learning in a Random World}, Springer.

\bibitem[Xie et al.(2024)]{xie2024finben}
Xie, Q., Han, W., Chen, Z., Xiang, R., Zhang, X., He, Y., Xiao, M., Li, D., Dai, Y., Feng, D., Xu, Y., Kang, H., Kuang, Z., Yuan, C., Yang, K., Luo, Z., Zhang, T., Liu, Z., Xiong, G., Deng, Z., Jiang, Y., Yao, Z., Li, H., Yu, Y., Hu, G., Huang, J., Liu, X.-Y., Lopez-Lira, A., Wang, B., Lai, Y., Wang, H., Peng, M., Ananiadou, S., and Huang, J., 2024, ``FinBen: A Holistic Financial Benchmark for Large Language Models'', in \emph{Advances in Neural Information Processing Systems}, Vol. 37.

\bibitem[Yan et al.(2024)]{yan2024crag}
Yan, S.-Q., Gu, J.-C., Zhu, Y., and Ling, Z.-H., 2024, ``Corrective Retrieval Augmented Generation'', arXiv:2401.15884.

\bibitem[Yao et al.(2023)]{yao2023react}
Yao, S., Zhao, J., Yu, D., Du, N., Shafran, I., Narasimhan, K., and Cao, Y., 2023, ``ReAct: Synergizing Reasoning and Acting in Language Models'', in \emph{International Conference on Learning Representations}.

\bibitem[Zhao et al.(2024)]{zhao2024findver}
Zhao, Y., Long, Y., Jiang, T., Wang, C., Chen, W., Liu, H., Tang, X., Zhang, Y., Zhao, C., and Cohan, A., 2024, ``FinDVer: Explainable Claim Verification over Long and Hybrid-Content Financial Documents'', in \emph{Proceedings of the 2024 Conference on Empirical Methods in Natural Language Processing}, pp. 14739--14752.

\bibitem[Zhu et al.(2021)]{zhu2021tatqa}
Zhu, F., Lei, W., Huang, Y., Wang, C., Zhang, S., Lv, J., Feng, F., and Chua, T.-S., 2021, ``TAT-QA: A Question Answering Benchmark on a Hybrid of Tabular and Textual Content in Finance'', in \emph{Proceedings of the 59th Annual Meeting of the Association for Computational Linguistics}, pp. 3277--3287.


\bibitem[Guan et al.(2026)]{guan2026finverify}
Guan, M., Chen, Q., Meng, G., and Liu, X.-Y., 2026, ``FinVerify: A Verification Framework for Financial Agents'', arXiv:2611.20423.

\bibitem[He et al.(2026)]{he2026chrono}
He, S., Lv, Y., Manela, A., and Wu, J., 2026, ``Chronologically Consistent Large Language Models'', \emph{Journal of Financial Economics}, Vol. 185, 104374.

\bibitem[Joren et al.(2025)]{joren2025sufficient}
Joren, H., Zhang, J., Ferng, C.-S., Juan, D.-C., Taly, A., and Rashtchian, C., 2025, ``Sufficient Context: A New Lens on Retrieval Augmented Generation Systems'', in \emph{International Conference on Learning Representations}.

\bibitem[Karger et al.(2025)]{karger2025forecastbench}
Karger, E., Bastani, H., Yueh-Han, C., Jacobs, Z., Halawi, D., Zhang, F., and Tetlock, P. E., 2025, ``ForecastBench: A Dynamic Benchmark of AI Forecasting Capabilities'', in \emph{International Conference on Learning Representations}.

\bibitem[Mao et al.(2025)]{mao2025selectllm}
Mao, Y., Durand, A., Mehrasa, N., He, Y., and Ester, M., 2025, ``Calibrating LLMs for Selective Prediction: Balancing Coverage and Risk'', in \emph{Advances in Neural Information Processing Systems}, Vol. 38.

\bibitem[Tao et al.(2026)]{tao2026tempfinrag}
Tao, L., Li, Z., Ji, Y., Liao, C.-T., Xiao, X., and Cao, X., 2026, ``TempFinRAG: Multimodal Temporal Retrieval-Augmented Generation for Point-in-Time Financial Question Answering'', \emph{Symmetry}, Vol. 18, No. 9, 1498.

\bibitem[Yin et al.(2026)]{yin2026epsbench}
Yin, Y., Xu, C., Ye, L., Peng, J., and Zhang, C., 2026, ``EPSBench: Benchmarking LLM Agents for Earnings Forecasting'', arXiv:2610.38293.

\bibitem[Torres et al.(2026)]{torres2026finabstain}
Torres, D., Cheng, W., and Huang, H., 2026, ``FinAbstain: Uncertainty-Calibrated Multimodal RAG for Selective Financial Forecasting'', arXiv:2607.24875.

\bibitem[Wang et al.(2026)]{wang2026ane}
Wang, Z., Zeng, Y., Gong, Z., Guo, Y., Zhu, F., Zhang, H., Zhang, W., and Zuo, W., 2026, ``AnE: Pushing the Reasoning Frontier of Multimodal LLMs via Anchor Evolution'', in \emph{European Conference on Computer Vision (ECCV)}, LNCS Vol. 17016, pp. 410--428.

\end{thebibliography}
}
